\section{1. Introduction}

Light field imaging captures both spatial and angular information of light rays, providing rich geometric cues for accurate depth estimation [1]. This capability significantly advances downstream applications in autonomous navigation, robotic manipulation, and immersive virtual reality [2]. By analyzing epipolar plane image (EPI)s and defocus cues, modern algorithms achieve high precision in controlled environments. Consequently, light field depth estimation remains a central research topic in computational photography and 3D computer vision.

Despite these advances, accurately recovering depth in real-world scenes remains a substantial challenge due to the prevalence of non-Lambertian materials [3]. Most existing methods rely on the Lambertian assumption, enforcing photo-consistency $I(x, y, u, v) \approx I(x, y, u', v')$ across different viewpoints. However, real-world objects frequently exhibit complex reflectance properties, including specular reflections, anisotropic scattering, and mixed material compositions. These interactions violate the underlying photo-consistency constraint, introducing misleading bidirectional reflectance distribution function (BRDF) cues that severely corrupt geometric reasoning.

The severity of this issue becomes evident when evaluating state-of-the-art models on diverse material subsets. For instance, a strong baseline achieves a mean absolute error (MAE) of 0.081 on mixed urban scenes, demonstrating robust performance under partial Lambertian conditions. In stark contrast, its MAE deteriorates to 0.411 on strictly non-Lambertian subsets, significantly exceeding the acceptable threshold of 0.25. More critically, visual inspections reveal that predictions in specular regions completely degenerate into mere replications of input textures. The network entirely fails to extract structural depth, outputting flat surfaces instead of true geometric structures.

This severe degradation highlights the inadequacy of purely data-driven approaches that ignore underlying physical optics. When specular highlights shift across viewpoints, traditional epipolar slopes $\frac{du}{dx}$ become distorted, rendering standard geometric priors invalid. Therefore, developing a unified depth estimation framework that explicitly models light-matter interactions is highly necessary. Without incorporating physical reflectance models, light field systems will remain confined to simplified laboratory settings. Addressing this non-Lambertian bottleneck is necessary for deploying robust 3D perception in unconstrained, real-world environments.

Traditional light field depth estimation methods primarily rely on epipolar plane image (EPI) analysis and focal stack evaluation [4]. Methodologically, these approaches strictly enforce the Lambertian assumption, presuming constant radiance across varying viewpoints. However, this idealized physical premise fails in real-world environments. Technically, specular reflections and anisotropic scattering severely distort epipolar line structures, violating the photo-consistency constraint. Consequently, slope extraction algorithms produce erroneous disparity values in non-Lambertian regions. Furthermore, even in purely diffuse areas, complex textures often cause slope ambiguity, limiting the accuracy of conventional geometric reasoning [5].

To mitigate these geometric ambiguities, early deep learning methods formulate depth estimation as an end-to-end regression task [6]. While convolutional neural networks effectively capture local spatial-angular correlations, they inherently lack explicit physical priors. This methodological flaw forces models to rely on statistical shortcuts rather than genuine geometric cues. Specifically, when processing non-Lambertian surfaces, these networks frequently memorize high-frequency specular highlights. As a result, the depth predictions degenerate into mere replications of input textures, yielding substantial errors. Additionally, the black-box optimization process makes these models highly susceptible to overfitting, particularly when trained on limited non-Lambertian samples .

More recently, advanced architectures have incorporated attention mechanisms and multi-scale feature fusion to improve robustness in mixed scenes [7]. Despite these architectural enhancements, their underlying physical modeling remains superficial. Most contemporary methods employ heuristic post-processing or simple binary masks to separate diffuse and specular components. This simplistic treatment fails to capture the continuous variations of surface roughness and complex scattering phenomena. Technically, the absence of fine-grained bidirectional reflectance distribution function constraints leads to severe depth artifacts at material boundaries. Moreover, the massive parameter spaces of these networks exacerbate training instability, often trapping the optimization in local minima without guaranteeing physical consistency [8].

The aforementioned limitations highlight a significant gap in current research. Existing techniques either depend on overly idealized physical assumptions or employ physically ungrounded data-driven mappings. Therefore, there is an urgent need for a unified approach that explicitly models light-matter interactions at the physical optics level. Such a method must integrate rigorous physical priors with component-aware depth estimation strategies to accurately resolve complex non-Lambertian and mixed scenes.

In this paper, we propose a unified dual-mask physical model to address non-Lambertian light field depth estimation. Our approach explicitly integrates physical reflection priors with deep feature extraction to overcome photo-consistency failures in complex scenes. Specifically, we employ a dual-mask generation module that predicts spatial occlusion and angular material masks. These masks subsequently guide the cost volume aggregation to filter out specular highlights and occluded regions. Furthermore, we introduce a component-aware depth estimation strategy that adaptively fuses disparity hypotheses using material-specific physical models.

The main contributions of this paper are summarized as follows:

\begin{itemize}
  \item \textbf{Contribution 1}: We propose a dual-mask physical modeling approach based on MRI-inspired angular frequency analysis. By applying a 2D discrete Fourier transform to the $9 \times 9$ angular sampling matrix, we analyze spectral energy distributions to classify reflectance types. This mechanism generates a medium mask for surface roughness and an angular direction mask for wavevector variations. Consequently, it provides reliable physical priors and high-precision material distribution maps for complex mixed scenes.
  \item \textbf{Contribution 2}: We develop a component-aware depth estimation strategy driven by inverse wavevector analysis to resolve texture replication issues. This strategy implements a three-tier progressive mechanism, ranging from lightweight spectral classification to full wavevector-level physical parsing. It adaptively applies epipolar plane image (EPI) methods for diffuse regions while utilizing photometric stereo for specular areas. Finally, it fuses component depths via pixel-level confidence weighting, substantially improving robustness in non-Lambertian environments.
  \item \textbf{Contribution 3}: We establish a physics-driven, three-stage progressive validation and training framework to ensure architectural interpretability. This framework enforces strict phase gates that block subsequent model training until prior physical hypotheses pass quantitative verification. By replacing single end-to-end loss functions with physics-consistency metrics, this approach prevents overfitting during black-box optimization. Therefore, it guarantees physical interpretability and enhances generalization ability on limited non-Lambertian training samples.
\end{itemize}
Extensive experiments on four public light field datasets demonstrate the superiority of our proposed method. The model achieves an overall mean absolute error of less than 0.20, outperforming existing baseline networks. Notably, it yields considerable improvements in non-Lambertian and mixed subsets, validating the effectiveness of unified physical modeling.